%% ma: L10-B15-0102
%% lop: 10
%% chuong: 6
%% bai: 15
%% dang: 10.15.4
%% loai: TN4
%% muc: TH
%% dap_an: "B"
%% nguon: "Ảnh giáo viên gửi 25/11/2026 (lô 2): Bài tập trắc nghiệm – Toán 10 Bài 15 Hàm số (sự biến thiên, đồ thị), Phần I Câu 13"
%% kien_thuc: "Lớp 10 Bài 15 (đồng biến, nghịch biến); hàm số bậc nhất y=ax+b đồng biến khi a>0 (lớp 9)"
%% trang_thai: chinh-thuc
\begin{ex}
Cho hàm số $f(x)=4-3x$. Khẳng định nào sau đây đúng?
\choice{Hàm số đồng biến trên $\left(-\infty;\dfrac43\right)$.}{\True Hàm số nghịch biến trên $\left(\dfrac43;+\infty\right)$.}{Hàm số đồng biến trên $\mathbb R$.}{Hàm số đồng biến trên $\left(\dfrac34;+\infty\right)$.}
\loigiai{$a=-3<0$ nên hàm số nghịch biến trên $\mathbb R$, do đó nghịch biến trên $\left(\tfrac43;+\infty\right)$.}
\end{ex}
